\scanpage{107}
\begin{center}- \mspage{168} -\end{center}
\margin{IX}
13.1.4. Bien entendu, la suite exacte (13.1.3) est \uline{fonctorielle en} $A$, et commute à tout changement de base $S'\longrightarrow S$. De cette façon on trouve une équivalence de catégories, compatible aux changements de base, entre la catégorie des groupes analytiques relatifs sur $S$ lisses et à fibres connexes, et la catégorie des couples $(E,R)$ d'un fibré vectoriel localement libre $E$ sur $S$ et d'un sous-groupe analytique $R$ de $E$ étale sur $S$.

13.1.5. Lorsque $A$ est un groupe ``\uline{abéloïde}'' sur $S$, par quoi nous entendons qu'il est lisse, propre et séparé sur $S$ et à fibres connexes, de sorte qu'on a en particulier $A=A^0$, alors $R$ est un groupe localement constant sur $S$, dont le rang en chaque point $s$ (en tant que $\mathbb Z$-module) est égal à $2\operatorname{rang}_s(E)$, et réciproquement. De cette façon on trouve l'interprétation bien connue d'un espace abéloïde sur $S$ en termes d'un ``\uline{système local}'' $R$ sur $S$ et d'une structure de fibré vectoriel analytique \uline{complexe} sur le fibré analytique réel $R\otimes_{\mathbb Z}\mathbb R$ défini par $R$.

13.2. \uline{Complétion formelle le long d'une fibre}. Nous supposons donné désormais un point $s\in S$, et allons travailler localement au voisinage de $s$. Pour tout entier $n\geq0$, désignons par $S_n$ le n.ième voisinage infinitésimal de $s$ dans $S$, et pour tout espace analytique $X$ sur $S$, soit de même $X_n=X\times_S S_n$ ; pour $n$ variable, le système des $X_n$ forme alors un système inductif d'espaces analytiques sur $S$, qu'on dénotera par $\widehat X$ et appellera le \uline{complété formel de} $X$ \uline{le long de} $X_s$. Il dépend évidemment fonctoriellement de $X$, et sa formation commute aux limites projectives

\margin{486}
\margin{486}
\margin{106}
